% ---------------------------------------------------------------------------
% Section: Phase 11 -- a Codex-style agent harness for small models on a CPU
% ---------------------------------------------------------------------------
\section{Phase 11: A Codex-Style Agent Harness for Small Models on a CPU}
\label{sec:harness2}

Phase~10 took leaves of Codex (the patch format, output truncation, the sandbox
parameters) and kept a simple agent loop. Its live run (Section~\ref{sec:harness})
showed where that loop loses time and accuracy on this machine: every step
re-read almost its whole prompt, three of five actions were malformed, and one
model call overran the run's cap. Phase~11 (25--26 October 2026) ports the
\emph{harness} mechanisms of Codex's agent core \cite{openaicodex2026} (its turn
loop, history, compaction, tool grammar and budget handling), adapted to a
chat-completions server, a 1.7B and a 4B model, and four CPU cores. The source
survey is in \texttt{docs/research/harness\_integration.md} (Section~5) and the
decision in \texttt{docs/decisions/006-codex-harness.md}.

\subsection{Where an agent step spends its time}

A step's wall time is, to first order,
\begin{equation}
T_{\text{step}} \;=\; \frac{n_{\text{prompt}}-n_{\text{cached}}}{r_p} \;+\; \frac{n_{\text{dec}}}{r_d} \;+\; T_{\text{tools}},
\label{eq:stepcost}
\end{equation}
where $n_{\text{cached}}$ prompt tokens are served from the server's cache, $r_p$
is the prefill rate and $r_d$ the decode rate. On the reference machine System~2
prefills about $r_p=31$ tokens/s and decodes $r_d\approx6$--$8$ tokens/s
(measured 25 October 2026): a cold 1{,}969-token prompt took 64\,s to prefill, and
the same prompt with one more exchange appended took 3\,s, because only the 21 new
tokens were computed. Prefill is therefore not negligible next to decoding, and
whether the cache is hit decides it.

\subsection{What Codex does, and what carries over}

Codex speaks only the Responses API and assumes frontier models, so nothing is
embedded; Table~\ref{tab:codexport} lists the mechanisms and their adaptation.

\begin{table*}[t]
\centering
\caption{Codex agent-harness mechanisms (\texttt{codex-rs} at \texttt{970111f6}) and their adaptation in \Sys{} (crate \texttt{nafs-harness}).}
\label{tab:codexport}
\scriptsize
\begin{tabular}{@{}p{0.2\textwidth}p{0.36\textwidth}p{0.37\textwidth}@{}}
\toprule
Mechanism & In Codex & In \Sys{} \\
\midrule
Append-only history & \texttt{context\_manager}: items are only appended; a change of state is a new item; requests are sent incrementally when the new input extends the previous one & \texttt{context.rs}: the transcript only grows; the task message is never rewritten, so llama-server's slot cache reuses every earlier exchange (Proposition~\ref{prop:prefix}) \\
Compaction & at 90\,\% of the window, a model-written handoff summary plus the user messages ($\le$20k tokens) replace the history & at \texttt{compact\_at\_tokens}, a \emph{harness-written} state replaces it: the task, the current files, the last verdict and the attempts so far; same prefix as the first message, so the system prompt and task stay cached \\
Freeform \texttt{apply\_patch} & the patch is a freeform tool constrained by a Lark grammar & \texttt{grammar.rs}: the grammar in GBNF, rebuilt per step (only existing files can be updated, no unified-diff ranges, the step's allowed actions), with the reasoning span kept free \\
Dispatch at item done & a tool call runs when its output item is complete, while the response still streams & the stream ends when the action is complete (unconstrained requests); under the grammar the server itself stops there \\
Truncation on insertion & tool output capped (10k tokens) as it enters history & observations and evidence middle-truncated to 2\,kB, files to 6\,kB (a 10k-token output would take minutes to prefill here) \\
Budget reminders & \texttt{rollout\_budget}, \texttt{token\_budget}: remaining-budget notes appended to history & ``Steps left: $n$'' appended to each failing observation; the last step offers no \texttt{run} \\
Turn diff & \texttt{turn\_diff\_tracker}: the turn's net change as a unified diff & the run's net diff returned with the files \\
Cancellation & a turn aborts in flight & the run's wall-clock cap holds \emph{inside} a model call (the stream is dropped, which stops the server) \\
\midrule
Not taken & parallel tool calls under a read/write lock; websocket reuse; remote and async compaction; MCP; approvals & one action per reply (the verifier runs after every change, so a second action would duplicate it); the rest needs server state llama.cpp does not have \\
\bottomrule
\end{tabular}
\end{table*}

\subsection{An append-only transcript}

\begin{proposition}[Prefix reuse]\label{prop:prefix}
Let a server reuse the longest common prefix between a request and the token
sequence cached in its slot, and let the slot hold the previous request of the
run. In the Phase~10 layout, the files are part of the first user message, so
after a step that changed a file the reusable prefix ends at the first changed
byte of the file view: each such step prefills at least the rest of its prompt,
and a run of $K$ steps prefills $\Theta(K\cdot L)$ tokens for prompts of length
$L$. In the append-only layout the reusable prefix covers the whole previous
request, so step $k$ prefills only the previous reply and its observation, and a
run prefills each transcript token once: $\Theta(L_K)$ for a final transcript of
length $L_K$, plus one task message per compaction.
\end{proposition}
\begin{proof}[Sketch]
In the append layout request $k$ is request $k-1$ followed by the assistant
message and the observation of step $k-1$. The cached sequence is request $k-1$
followed by the generated tokens; the chat template renders the earlier assistant
turn without its reasoning, so the common prefix is exactly request $k-1$. In the
reset layout the first divergence is inside the first user message whenever a
file changed. Summing over steps gives the two totals.
\end{proof}

Codex keeps its history append-only for the same reason (its incremental requests
fail over to full ones on any mismatch), and prefix-sharing caches such as
SGLang's RadixAttention rest on the same invariant \cite{zheng2024sglang}. One consequence is that the model must
learn the new state from the observation rather than from a rewritten task
message: after a patch the observation carries the edited files, and after a
failing verdict the verifier's evidence. When the estimated transcript exceeds
\texttt{compact\_at\_tokens} (2{,}600 tokens by default, which leaves room for a
reply in System~1's 4{,}096-token slots) it is compacted into the task message
with an ``already tried'' list in place of Codex's model-written summary: at
6--8 tokens/s a summary would cost minutes, and the facts it needs are already
structured.

\subsection{Actions under a grammar}

Codex constrains \texttt{apply\_patch} with a grammar; llama-server accepts a GBNF
grammar in the request. The port specialises the grammar to each step: the
\texttt{Update} and \texttt{Delete} targets are the workspace's files, an
\texttt{@@} header cannot start with a range, a \texttt{run} block can name only
a Python file that exists, and the last step offers only edits. When the grammar
is complete only end-of-generation can be sampled, so nothing after the action is
decoded. \texttt{SUBMIT} is no longer offered: the verifier ends the run, so
\texttt{SUBMIT} only spent a step.

A reasoning model needs one more rule. llama-server applies a request grammar
from the first sampled token; with a plain grammar Qwen3-4B-Thinking was forced to
emit the answer in place of its reasoning (measured on llama.cpp \texttt{797de63}).
The grammar therefore opens with the reasoning span, any text up to
\texttt{</think>}:
\begin{lstlisting}
root  ::= think action
think ::= ( [^<] | "<" [^/] | "</" [^t] | "</t" [^h] | "</th" [^i]
          | "</thi" [^n] | "</thin" [^k] | "</think" [^>] )* "</think>" [ \n]*
action ::= patch | whole | run
whole ::= "```python\n" codeline* "```"
ctx   ::= "@@" ( " " [^-\n] [^\n]* )? "\n"
\end{lstlisting}
With this grammar and a 60-token reasoning budget the server closed the reasoning
at the budget and sampled a well-formed answer. Reasoning stays unconstrained,
which matters because format restrictions on the reasoning itself can cost
accuracy \cite{tam2024speakfreely}; grammar-constrained decoding is otherwise
cheap on the server side \cite{willard2023outlines,dong2024xgrammar}.

\subsection{Which edits to offer}

Codex edits only through patches. Small models are worse at diff formats than at
whole-file rewrites \cite{diffxyz2025,aider2024editformats}, and a patch saves
answer tokens only when the file is long (Eq.~\ref{eq:patch}). The live runs
agree: in Phase~10 Qwen3-4B wrote a unified-diff header, an update of a file that
did not exist, and an update hunk without line prefixes; in the first Phase~11
smoke run Qwen3-1.7B wrote well-formed patches under the grammar that appended a
second copy of the file. \texttt{edit\_format = "auto"} therefore offers the
whole \texttt{solution.py} while it is missing or at most 80 lines long and a
patch or the whole file above that; \texttt{"any"} offers both at every step, as
Phase~10 did. Patches remain the edit for long files, where Eq.~\ref{eq:patch}
favours them.

\subsection{Smaller pieces}

A workspace the verifier already judged is not judged again (a model that
restores an earlier file is told so). Tracebacks name files relative to the run
(\texttt{File "runner.py"}), not the sandbox's random directory. Each step
records the server's token accounting (prompt tokens computed and served from the
cache, decode tokens, finish reason), and each run returns the totals and its net
diff. The Phase~10 loop remains selectable (\texttt{context = "reset"}), and the
request may override \texttt{context}, \texttt{constrain}, \texttt{stream\_cut}
and \texttt{edit\_format}, which is how the A/B below was run.

\subsection{Results}

% Generated by scripts/gen_harness_ab.py from benchmarks/harness_ab*.json; do not edit by hand.
\begin{table*}[t]
\centering
\caption{Code agent, interleaved live A/B on the reference machine (\texttt{benchmarks/harness\_ab.json}). Same tasks, model and budget (6 steps) per arm; arm order rotated per task. Malformed: steps whose reply held no action or a patch that did not apply. Prefill: prompt tokens the server computed; cache: share of prompt tokens served from its cache. Wall times are from a shared host and are not speed claims.}
\label{tab:harnessab}
\small
\begin{tabular}{@{}lrrrrrrr@{}}
\toprule
Arm & Verified & Calls & Malformed & Prefill tok. & Cache & Decoded tok. & Wall (s) \\
\midrule
\multicolumn{8}{@{}l}{\emph{System~1 (Qwen3-1.7B), 10 tasks}} \\
Phase 10 loop & 2/10 & 50 & 23/50 & 6{,}415 & 75\,\% & 2{,}793 & 356 \\
Phase 11 (auto edits) & 4/10 & 41 & 1/41 & 7{,}345 & 78\,\% & 4{,}146 & 504 \\
Phase 11, patch always & 0/10 & 60 & 16/60 & 15{,}393 & 76\,\% & 10{,}706 & 1{,}468 \\
\midrule
\multicolumn{8}{@{}l}{\emph{System~2 (Qwen3-4B-Thinking), 3 tasks}} \\
Phase 10 loop & 0/3 & 15 & 13/15 & 1{,}449 & 85\,\% & 9{,}907 & 2{,}695 \\
Phase 11 (auto edits) & 3/3 & 3 & 0/3 & 741 & 23\,\% & 1{,}818 & 451 \\
\bottomrule
\end{tabular}
\end{table*}
\begin{table*}[t]
\centering
\caption{Round 2, on the build with value-carrying evidence (\texttt{benchmarks/harness\_ab\_round2.json}): \texttt{unstick} on System~1 and the reasoning budget per step on System~2. Columns as in Table~\ref{tab:harnessab}.}
\label{tab:harnessab2}
\small
\begin{tabular}{@{}lrrrrrrr@{}}
\toprule
Arm & Verified & Calls & Malformed & Prefill tok. & Cache & Decoded tok. & Wall (s) \\
\midrule
\multicolumn{8}{@{}l}{\emph{System~1 (Qwen3-1.7B), 10 tasks}} \\
Phase 11 without unstick & 5/10 & 35 & 0/35 & 6{,}647 & 79\,\% & 3{,}926 & 523 \\
Phase 11 with unstick & 7/10 & 27 & 0/27 & 3{,}283 & 80\,\% & 2{,}575 & 305 \\
\midrule
\multicolumn{8}{@{}l}{\emph{System~2 (Qwen3-4B-Thinking), 6 tasks}} \\
Phase 11, 512 reasoning tokens & 6/6 & 7 & 0/7 & 276 & 89\,\% & 4{,}175 & 718 \\
Phase 11, 256 reasoning tokens & 6/6 & 14 & 5/14 & 1{,}612 & 81\,\% & 4{,}748 & 858 \\
\bottomrule
\end{tabular}
\end{table*}
\begin{table*}[t]
\centering
\caption{Round 3, on the build whose grammar keeps patches out of code blocks (\texttt{benchmarks/harness\_ab\_round3.json}): the System~2 reasoning budget fixed or raised after a failure. Columns as in Table~\ref{tab:harnessab}.}
\label{tab:harnessab3}
\small
\begin{tabular}{@{}lrrrrrrr@{}}
\toprule
Arm & Verified & Calls & Malformed & Prefill tok. & Cache & Decoded tok. & Wall (s) \\
\midrule
\multicolumn{8}{@{}l}{\emph{System~2 (Qwen3-4B-Thinking), 6 tasks}} \\
512 reasoning tokens every step & 6/6 & 6 & 0/6 & 140 & 93\,\% & 3{,}499 & 560 \\
256 on the first step, then 512 & 6/6 & 6 & 0/6 & 6 & 100\,\% & 1{,}961 & 282 \\
\bottomrule
\end{tabular}
\end{table*}

% Phase 11 results text. Numbers come from harness_ab_gen.tex (scripts/gen_harness_ab.py).
\textbf{Design of the A/B.} \texttt{scripts/harness\_ab.py} runs ten small Python
tasks (six fixes of a buggy \texttt{solution.py}, four written from an empty
workspace; four or five hidden asserts each) through the running gateway, once
per arm, with the same model and a budget of six steps; the order of the arms is
rotated from task to task so that a drift in machine load does not favour one
arm. The arms are the Phase~10 loop (context rebuilt every step, no grammar, no
stream cut, patches always offered), Phase~11 with \texttt{edit\_format =
"auto"}, and, on System~1, Phase~11 with a patch offered at every step, which
isolates the edit policy. The host was shared with the owner's other model
servers (load average 5--9), so wall times are reported but are not speed
claims; token counts and outcomes do not depend on load.

\textbf{System~1.} Table~\ref{tab:harnessab} gives the totals. The Phase~11 arm
verified \HABOneNewVerified{} of \HABOneNewN{} tasks in \HABOneNewCalls{} model
calls; the Phase~10 loop verified \HABOneOldVerified{} in \HABOneOldCalls{}. The
difference in verified tasks lies on two discordant tasks (Phase~11 verified the
palindrome fix and \texttt{flatten}, where Phase~10 did not; no task went the
other way), so with ten tasks it is indicative and not significant (exact
McNemar $p=0.5$). The mechanism behind it is not in doubt: \HABOneOldMalformed{}
of the Phase~10 loop's \HABOneOldActions{} steps held no action it could parse
(the 1.7B model answered in prose or with \texttt{SUBMIT}), against
\HABOneNewMalformed{} of \HABOneNewActions{} under the grammar, a reply that ran
out of tokens before its fence closed. Offering a patch at every step was the
worst arm: \HABOneAnyVerified{} of \HABOneAnyN{} verified, \HABOneAnyMalformed{}
of \HABOneAnyActions{} steps malformed, and the most tokens decoded
(\HABOneAnyDecoded{}). Most of its patches \emph{applied} and still broke the
file: the 1.7B model prefixed added lines with extra plus signs (\texttt{++ def
rle(s):}), left a mangled \texttt{End Patch} line inside the file, or replaced
the whole file with \texttt{python check.py}; the grammar fixes the syntax of a
patch, not its sense. Prefill was not where System~1 lost time: llama-server's chunked cache
reuse (\texttt{--cache-reuse 256}) and its host-memory prompt cache served
\HABOneOldHit\,\% of the Phase~10 loop's prompt tokens as well, against
\HABOneNewHit\,\% for the append-only transcript. Phase~11 decoded more
(\HABOneNewDecoded{} against \HABOneOldDecoded{} tokens) because its steps are
whole files rather than unparseable prose, and spent \HABOneNewWall\,s against
\HABOneOldWall\,s in all: \HABOneNewWallPerVerified\,s per verified task against
\HABOneOldWallPerVerified\,s.

% Phase 11 results, System 2 (round 1). Numbers from harness_ab_gen.tex.
\textbf{System~2.} System~2 is the code agent's default model, and three tasks
were run per arm (a System~2 step takes minutes on this host). The Phase~11 arm
verified all \HABTwoNewVerified{} in \HABTwoNewCalls{} model calls, one per task,
in \HABTwoNewWall\,s in all; the Phase~10 loop verified
\HABTwoOldVerified{} of \HABTwoOldN{} in \HABTwoOldCalls{} calls and
\HABTwoOldWall\,s. Two of its runs ended at the 900\,s cap, which now holds
inside a model call (both ended at 900.0\,s; in Phase~10 one call had run to
1{,}041\,s), and the third spent its six steps. \HABTwoOldMalformed{} of its
\HABTwoOldActions{} steps were malformed: update hunks whose lines lacked the
\texttt{+}/\texttt{-}/space prefix or the \texttt{@@} marker, and replies with no
action at all, two of them cut off by the token limit. Under the grammar no step
was malformed. Two of the three Phase~11 steps are worth a note: the files were
short, so \texttt{"auto"} offered only a whole file, and the model wrote a Codex
patch \emph{inside} the \verb|```python| block, which the grammar allows (a code
line may start with \texttt{***}); the parser applies a patch envelope wherever it
appears, both applied, and both runs verified. Every Phase~11 call spent its full
reasoning budget (511 tokens): decoding took \HABTwoNewDecodeS\,s of the
\HABTwoNewWall\,s and prefill \HABTwoNewPrefillS\,s, which is why round~2
measures the budget. With three tasks this is a demonstration rather than an
estimate, but the direction agrees with System~1 and the mechanism is visible in
every step.

% Phase 11 results, rounds 2 and 3. Numbers from harness_ab_macros.tex.
\textbf{Round 2: breaking repetition, and the reasoning budget.} Round~1 showed
a failure the grammar cannot touch: a small model that rewrites the same failing
file step after step, anchored on its own previous reply (System~1 reproduced
an already-judged file in four consecutive steps of one run). Two changes went
into the build for round~2. Failing asserts now carry the value they got
(\texttt{assert f(x) == y -> AssertionError: got 5}), the concise execution
feedback that self-debugging relies on \cite{chen2024selfdebug}; and
\texttt{unstick}: after a step that reproduced files the verifier had already
judged, the transcript is compacted before the next call (the model no longer
sees, and copies, its repeated reply; the attempts list says what was tried) and
the step is sampled 0.3 hotter per repeat. Table~\ref{tab:harnessab2} gives the
interleaved comparison on the same ten System~1 tasks: with \texttt{unstick} the
agent verified \HABROneUVerified{} in \HABROneUCalls{} calls and \HABROneUWall\,s,
without it \HABROneNoUVerified{} in \HABROneNoUCalls{} calls and
\HABROneNoUWall\,s; the two discordant tasks both favour \texttt{unstick}, so it is
on by default. The value-carrying evidence was in both arms and is not isolated;
the Phase~11 arm verified 4 of 10 in round~1 and \HABROneNoUVerified{} without
\texttt{unstick} in round~2, which is consistent with it helping but is not an
interleaved comparison.

On System~2 every call of round~1 spent its full reasoning budget, so round~2
compared 512 and 256 reasoning tokens per step on six tasks. Both verified all
six. Halving the budget halved the time of the four tasks it solved in one call
(about 55\,s against 100\,s), but two tasks then needed four and six calls, \HABRTwoSmallMalformed{}
of its \HABRTwoSmallActions{} steps were malformed, and the arm took
\HABRTwoSmallWall\,s against \HABRTwoBigWall\,s: a smaller budget is not faster
overall. The malformed steps exposed a gap in the grammar. With a short file only
a whole file is offered, and System~2 wrote a Codex patch \emph{inside} the
\verb|```python| block, which the grammar allowed (round~1's two such patches
applied; round~2's four did not). A code line may no longer start with
\texttt{*} (a top-level Python line never does), so the edit policy is now
enforced by the grammar, not only by the prompt.

\input{sections/harness2_results_r3.tex}



